\section{Problem Formulation}\label{sec:problem}
\paragraph*{Plant (the digital twin).} A brushed DC motor with armature resistance $R$, torque and back-EMF constants $K_t=K_e=K$ (ideal motor), rotor inertia $J_m$ and viscous friction $B_m$ drives a gearbox of ratio $N=180$ and efficiency $\eta$ and a revolute joint with a thin-rod arm and a point-mass payload. Armature inductance is neglected in the base twin ($L=0$). With motor speed $\omega_m$ and joint angle $\theta$,
\begin{align}
V &= R\,i + K\,\omega_m, \qquad J_m^{\mathrm{tot}}\,\dot\omega_m = K\,i - B_m\,\omega_m - \tau_{\mathrm{load}}/(N\eta),\label{eq:plant}\\
\tau_{\mathrm{load}} &= J_{\mathrm{out}}\ddot\theta + m g\,\ell\,\sin\theta + \tau_{\mathrm{fric}},
\end{align}
with $\theta=0$ hanging, $J_m^{\mathrm{tot}}$ including the payload inertia reflected through the gearbox, and a joint friction term that is zero by default. The supply is \SI{12}{V}; the encoder resolves \num{9360} counts per output revolution; the loop runs at \SI{1}{kHz}. Parameters come from a datasheet with documented gaps and a stated inconsistency between the no-load speed and the stall torque, so all parameters carry a provenance tag and the twin is called preliminary.

\paragraph*{Controller.} A PID position controller with output saturation at the supply voltage, conditional anti-windup and a filtered derivative acts on the encoder error. Gains $(K_p,K_i,K_d)$ are the quantity to be chosen.

\paragraph*{Tuning problem and scenarios.} A scenario fixes the plant parameters (resistance, motor constant, rotor inertia, viscous friction, efficiency, drawn inside datasheet-derived ranges with one shared motor constant), a command delay of $0$--$2$ samples, encoder noise, the target angle in \SIrange{30}{180}{deg}, the payload mass (up to \SI{2}{kg}), the arm length, a gravity flag and the available voltage fraction. A scenario is \emph{infeasible} if the torque-feasibility margin (available stall torque over peak gravity torque on the path, an upper bound that ignores back-EMF) falls below one at any point, including after a mid-episode payload change.

\paragraph*{Episodes, return and metrics.} An episode lasts \SI{2}{s}; gains may change every \SI{20}{ms} (100 decisions). The \emph{return} is the sum over decisions of $-(w_e\,\bar e^2_t + w_u\,\bar V^2 + w_c\,\overline{|\Delta V|} + w_g\,\|\Delta\log g\|_1 + w_o\,\mathrm{os})$, with a time-weighted squared error normalised by the step size, mean squared voltage, voltage chatter, a gain-change penalty and overshoot; the weights are listed in Table~\ref{tab:hparams}. Return is the \textbf{primary metric}. The fraction of scenarios that end within \SI{2}{\%} of the target with overshoot at most \SI{10}{\%} is reported as a \textbf{descriptive} success rate only (it was not stable across seed sets, Section~\ref{sec:v3}).
